\section*{Reading notes}
\subsection*{\texttt{GELFAND-P001465}}
\par\noindent\leavevmode\phantomsection\label{GELFAND-0159-U001}\label{GELFAND-P001465:eq18-composition}\textbf{\texttt{GELFAND-0159-U001}} (eq18-composition).\nopagebreak[4]\par
In (18), a block BA is followed by joined blocks B and A, then a sum and two separate blocks. The primes on the upper j labels in the middle drawing resemble double primes; the last lower index beside A intersects a strand and resembles M, although m\_L stands below.\par
Candidates: Above B: j\_1'', j\_2'', j\_M; below A: j\_1, j\_2, j\_M; Some upper j labels have one prime; the last lower index beside A is L. Unresolved in the diplomatic source.\par
Historical source review: Original 2408×3480 pixels. Entire region (65,100,2040,975); contextual enlargements 3×: (65,100,715,735), (865,155,1365,975), (1370,220,2040,970); details 8×: (180,180,600,380), (930,185,1400,340), (1165,730,1320,865). Indices were not reconciled according to the expected composition rule. The description preserves the visible notation, including the M/L discrepancy. Reading the upper endings beside B is hindered by overlap with the heading.\par
Current source comparison: The overlapping heading still prevents a secure count of every prime. The lower M/L discrepancy remains literal. The partial diagram marks these labels explicitly as uncertain.\par\medskip
\par\noindent\leavevmode\phantomsection\label{GELFAND-0159-U002}\label{GELFAND-P001465:eq20-crossing}\textbf{\texttt{GELFAND-0159-U002}} (eq20-crossing).\nopagebreak[4]\par
At the lower-right end of the continuous diagonal in drawing (20), a short curved side stroke adjoins a pointed end. A second wing of an arrow cannot be separately distinguished. The labels and the gap in the other strand are legible.\par
Candidates: An arrow down and right; A retraced pointed end without an arrow. Unresolved in the diplomatic source.\par
Historical source review: Original 2408×3480 pixels; region (220,1250,515,1495), context (60,950,1810,1515). Examined in the original and at 3×, with a further independent enlarged comparison. The strand direction was not reconstructed from subsequent drawings. In (20′), there is no such side stroke at the end.\par
Current source comparison: The one-sided hook and pointed end remain visible at 5×/10×, but do not securely establish an arrow. The recreation preserves the crossing gap and omits an asserted direction at that end.\par\medskip
\par\noindent\leavevmode\phantomsection\label{GELFAND-0159-U003}\label{GELFAND-P001465:eq20-v-indices}\textbf{\texttt{GELFAND-0159-U003}} (eq20-v-indices).\nopagebreak[4]\par
The second index in R\textasciicircum{}\{j\_1 j\_2\} resembles 2. In all four V factors in the same line, the index beside j reads 1; a separate examination confirmed the same straight stroke with a lower foot.\par
Candidates: R\textasciicircum{}\{j\_1j\_2\}:V\textasciicircum{}\{j\_1\}\textbackslash{}otimes V\textasciicircum{}\{j\_1\}\textbackslash{}longrightarrow V\textasciicircum{}\{j\_1\}\textbackslash{}otimes V\textasciicircum{}\{j\_1\}. Unresolved in the diplomatic source.\par
Historical source review: Original 2408×3480 pixels; region (700,1220,1790,1400), 8×; general context (60,950,1810,1515), 3×. The repeated j\_1 was preserved literally. The neighboring drawings have colors j\_1,j\_2, but this does not justify substituting indices in the formula. Notational compatibility remains to be checked; no corrected version was asserted.\par
Current source comparison: All four lower indices remain visibly 1 at 3×/5×/10×. The separate tensor-space check does not replace them with the colors in the drawing.\par\medskip
\subsection*{\texttt{GELFAND-P001466}}
\par\noindent\leavevmode\phantomsection\label{GELFAND-0159-U004}\label{GELFAND-P001466:eq25-cancelled-exponent}\textbf{\texttt{GELFAND-0159-U004}} (eq25-cancelled-exponent).\nopagebreak[4]\par
In (25), the exponent on (-1) begins with m; to the right is a densely crossed-out sign with several fused strokes. The original index cannot reliably be separated from the cancellation. The exponent on q to the right reads −m\_1/2.\par
Candidates: m\_1; m with an illegible crossed-out sign to its right. Unresolved in the diplomatic source.\par
Historical source review: Original 2408×3480 pixels; region (1015,1265,1365,1375), 10×; context (105,795,1670,1540), 3×. The candidate m\_1 in the body is marked uncertain; the crossed-out sign was not separately recovered. No additional factor or numeral was added. The visible q exponent in (24) is m\_1 and in (25) is −m\_1/2; both were preserved without unification.\par
Current source comparison: The cancelled strokes remain inseparable even at 10×. The two different q exponents remain literal; a conditional cup/cap contraction is recorded separately with its conventions stated.\par\medskip
\par\noindent\leavevmode\phantomsection\label{GELFAND-0159-U005}\label{GELFAND-P001466:eq26-vertex-arrow}\textbf{\texttt{GELFAND-0159-U005}} (eq26-vertex-arrow).\nopagebreak[4]\par
In the right-hand drawing in (26), the arrow on the j\_2 branch below and to the right of the vertex is covered with dense repeated strokes. Short parts protrude on the upper-right and lower-left sides, permitting a direction toward the vertex.\par
Candidates: The j\_2 arrow points toward the vertex, up and left; The direction of the j\_2 arrow is illegible. Unresolved in the diplomatic source.\par
Historical source review: Original 2408×3480 pixels; entire region (140,1660,1365,2030); context (140,1550,1555,2465), 3×; detail (1190,1790,1320,1925), 10×. The legible branches and arrows on both sides of the equality are listed in SourceFigure. The ambiguous arrow direction was not set according to algebraic expectation.\par
Current source comparison: The heavily retraced mark remains directionally ambiguous at 10×. The partial recreation types j\_2 and traces its branch without choosing an arrowhead there.\par\medskip
\par\noindent\leavevmode\phantomsection\label{GELFAND-0159-U006}\label{GELFAND-P001466:eq28-coefficient-j}\textbf{\texttt{GELFAND-0159-U006}} (eq28-coefficient-j).\nopagebreak[4]\par
In (28), both outgoing lower strands are labeled j\_3. In the coefficient, the exponent of (-1) ends in −j, and the numerator of the q exponent begins j(j+1); a lower index 3 is not visible on those j symbols. The denominator of the q exponent is 4.\par
Candidates: (-1)\textasciicircum{}\{j\_1+j\_2-j\}q\textasciicircum{}\{(j(j+1)-j\_1(j\_1+1)-j\_2(j\_2+1))/4\}. Unresolved in the diplomatic source.\par
Historical source review: Original 2408×3480 pixels; region (140,2495,2040,3020), within the examined enlargement (140,2495,2040,3300), 3×. Literal j in the coefficient and j\_3 in the drawings were preserved; their possible identification was not inserted into the source as an established equality.\par
Current source comparison: The corrected overlapping 5×/10× details confirm the bare j, j(j+1), and denominator 4. The earlier mostly blank detail was inadequate evidence and has been explicitly superseded by those additional views.\par\medskip
\subsection*{\texttt{GELFAND-P001467}}
\par\noindent\leavevmode\phantomsection\label{GELFAND-0159-U201}\label{GELFAND-P001467:x58-y247-x646-y638}\textbf{\texttt{GELFAND-0159-U201}} (x58-y247-x646-y638).\nopagebreak[4]\par
Drawing (29): an upward arrow labeled t stands on the left, with a curved loop on the right. The upper and lower bends and two places in the upper middle portion are circled with ovals. The label m\_i stands above the left middle place, and j is on the right. Crossings, arrow-like strokes and gaps on the right are visible.\par
Candidates: none recorded. Unresolved in the diplomatic source.\par
Historical source review: Coordinates in the original 2408×3480 image: x=58–646,y=247–638, origin at the upper left. Checked on the whole image and at 3×. Overlapping strokes prevented establishing the exact passage of the branches at the crossings and every arrow direction. The visible drawing was described without inventing topology.\par
Current source comparison: The 5×/10× views retain crossing gaps and overlaid strokes. The partial recreation preserves the four ovals and clear labels; it does not certify the knot type or all hidden branch connections.\par\medskip
\par\noindent\leavevmode\phantomsection\label{GELFAND-0159-U202}\label{GELFAND-P001467:x656-y1355-x935-y1489}\textbf{\texttt{GELFAND-0159-U202}} (x656-y1355-x935-y1489).\nopagebreak[4]\par
In (30), between the product sign and the parentheses (\{m\}), is a small v-shaped sign. A pointer leads from it to the caption about the matrix elements calculated above.\par
Candidates: v; \textbackslash{}nu. Unresolved in the diplomatic source.\par
Historical source review: Coordinates in the original 2408×3480 image: x=656–935,y=1355–1489. Checked at 3× and 10×; the shape does not confidently distinguish Latin v from Greek nu. The formula gives the candidate v with an explicit uncertainty marker. No definition of this sign was found on the neighboring preceding sheet.\par
Current source comparison: The curved glyph remains compatible with both alphabets. Mathematical usage cannot resolve the handwriting.\par\medskip
\par\noindent\leavevmode\phantomsection\label{GELFAND-0159-U203}\label{GELFAND-P001467:x1935-y1670-x2000-y1726}\textbf{\texttt{GELFAND-0159-U203}} (x1935-y1670-x2000-y1726).\nopagebreak[4]\par
In (31), immediately after N\_-\textasciicircum{}i(D\_L), at the upper-right end of the q exponent, is a small crossed-out superscript sign.\par
Candidates: none recorded. Unresolved in the diplomatic source.\par
Historical source review: Coordinates in the original 2408×3480 image: x=1935–2000,y=1670–1726. Viewed within the formula at 3× and 5× and separately at 10×. The original sign beneath dense intersecting strokes is unreadable. The cancellation was retained descriptively.\par
Current source comparison: The intersecting strokes still prevent a reading. No exponent is supplied from an expected link-invariant formula.\par\medskip
\par\noindent\leavevmode\phantomsection\label{GELFAND-0159-U204}\label{GELFAND-P001467:x520-y2610-x734-y2730}\textbf{\texttt{GELFAND-0159-U204}} (x520-y2610-x734-y2730).\nopagebreak[4]\par
In the line “One can … construct knot inv. [source abbreviation]”, a short word between “can” and “construct” is crossed out.\par
Candidates: none recorded. Unresolved in the diplomatic source.\par
Historical source review: Coordinates in the original 2408×3480 image: x=520–734,y=2610–2730. The hidden word was not established after 3× and 10× examinations. The remainder of the line was transcribed.\par
Current source comparison: The word remains obscured by cancellation. The remainder of the sentence is translated independently.\par\medskip
\subsection*{\texttt{GELFAND-P001468}}
\par\noindent\leavevmode\phantomsection\label{GELFAND-0159-U205}\label{GELFAND-P001468:x1260-y33-x1400-y148}\textbf{\texttt{GELFAND-0159-U205}} (x1260-y33-x1400-y148).\nopagebreak[4]\par
In the top line, between the underlined surname “Ushveridze” and the name “Alexander”, is a short densely crossed-out entry.\par
Candidates: none recorded. Unresolved in the diplomatic source.\par
Historical source review: Coordinates in the original 2408×3480 image: x=1260–1400,y=33–148. Examined at 3× and 10×; no letters or initials beneath the strokes were read. The uncancelled given name and patronymic were preserved.\par
Current source comparison: The short cancellation remains unreadable. No initials are supplied from the uncancelled name.\par\medskip
\par\noindent\leavevmode\phantomsection\label{GELFAND-0159-U206}\label{GELFAND-P001468:x2040-y40-x2310-y150}\textbf{\texttt{GELFAND-0159-U206}} (x2040-y40-x2310-y150).\nopagebreak[4]\par
On the right in the top line, after “Grigorievich.”, is an entry crossed out with several nearly horizontal strokes.\par
Candidates: none recorded. Unresolved in the diplomatic source.\par
Historical source review: Coordinates in the original 2408×3480 image: x=2040–2310,y=40–150. Checked on the whole image and at 3× and 10×. The original signs merge with the cancelling strokes; no reading was established.\par
Current source comparison: The almost horizontal strokes continue to obscure the entry. The full uncancelled header is retained.\par\medskip
\par\noindent\leavevmode\phantomsection\label{GELFAND-0159-U207}\label{GELFAND-P001468:x1000-y1240-x1248-y1328}\textbf{\texttt{GELFAND-0159-U207}} (x1000-y1240-x1248-y1328).\nopagebreak[4]\par
Above the line after (4), “conditions on Δ …” is inserted. The final abbreviated word in the insertion begins with Cyrillic “нахо” and ends before a full stop.\par
Candidates: нахожд. [source abbreviation; finding]. Unresolved in the diplomatic source.\par
Historical source review: Coordinates in the original 2408×3480 image: x=1000–1248,y=1240–1328. Checked at 5× and 10×. The final joined letters are not entirely distinguishable; “нахожд.” was given as a candidate without expanding the abbreviation. The insertion and its position were marked in the text.\par
Current source comparison: The ending remains joined and uncertain. “Finding” in the body is an explicitly marked rendering of the candidate source abbreviation, not a claim that every abbreviated letter is secure.\par\medskip
\par\noindent\leavevmode\phantomsection\label{GELFAND-0159-U208}\label{GELFAND-P001468:x435-y2450-x570-y2578}\textbf{\texttt{GELFAND-0159-U208}} (x435-y2450-x570-y2578).\nopagebreak[4]\par
In (6), item 3), before the uncancelled equality Ω(λ)−Ω(−λ)=0, is a short crossed-out formula. A triangular sign resembling Δ can be distinguished on its left.\par
Candidates: none recorded. Unresolved in the diplomatic source.\par
Historical source review: Coordinates in the original 2408×3480 image: x=435–570,y=2450–2578. Checked at 10×. Bracket-like and letter strokes are covered by cancellation; the complete earlier fragment cannot be recovered.\par
Current source comparison: The triangular remnant does not establish the cancelled formula. The visible equality remains separate.\par\medskip
\par\noindent\leavevmode\phantomsection\label{GELFAND-0159-U209}\label{GELFAND-P001468:x1335-y2810-x1494-y2960}\textbf{\texttt{GELFAND-0159-U209}} (x1335-y2810-x1494-y2960).\nopagebreak[4]\par
To the right of the first pair of formulas in (7), after Ω(λ)=0;, a small fragment is crossed out. A triangular sign resembling Δ is distinguishable.\par
Candidates: none recorded. Unresolved in the diplomatic source.\par
Historical source review: Coordinates in the original 2408×3480 image: x=1335–1494,y=2810–2960. Viewed at 3× and 10×. The other signs beneath the dense strokes were not established; no mathematically expected expression was substituted.\par
Current source comparison: The remaining strokes still do not establish a formula beyond the triangular remnant.\par\medskip
\par\noindent\leavevmode\phantomsection\label{GELFAND-0159-U213}\label{GELFAND-P001468:x350-y1390-x1615-y1675}\textbf{\texttt{GELFAND-0159-U213}} (x350-y1390-x1615-y1675).\nopagebreak[4]\par
The first line of (5) is written sum\_\{i=1\}\textasciicircum{}M prime Delta(lambda\_i-lambda\_k) + sum\_\{alpha=1\}\textasciicircum{}N b\_alpha(lambda\_i-a\_alpha)=0. Beneath the first sum is “sum over i≠k”. Immediately after b\_alpha is an opening parenthesis, with no Delta; the index on lambda in the second parenthesis is specifically i.\par
Candidates: none recorded. Unresolved in the diplomatic source.\par
Historical source review: Coordinates in the original 2408×3480 image: x=350–1615,y=1390–1675. Checked on the whole image and at a 3× enlargement of formula (5). The absent Delta and the combination of a summed i with i in the second term require mathematical comparison with (1)–(4). This is a remark about a possible source error. All visible signs, indices and the annotation were preserved; no corrected equation was substituted.\par
Current source comparison: The corrected 5×/10× details confirm the omitted Δ and literal λ\_i in the second summand. Some earlier detail boxes clipped the formula; the additional views explicitly remedy that limitation. The separate residue calculation does not emend (5).\par\medskip
\subsection*{\texttt{GELFAND-P001469}}
\par\noindent\leavevmode\phantomsection\label{GELFAND-0159-U210}\label{GELFAND-P001469:x825-y1238-x1030-y1340}\textbf{\texttt{GELFAND-0159-U210}} (x825-y1238-x1030-y1340).\nopagebreak[4]\par
In the insertion “in the case … (7)\_1” above item В, a crossed-out entry precedes the uncancelled reference (7)\_1.\par
Candidates: none recorded. Unresolved in the diplomatic source.\par
Historical source review: Coordinates in the original 2408×3480 image: x=825–1030,y=1238–1340. Viewed at 3× and 10×. The earlier number or formula sign cannot be distinguished beneath intersecting strokes. The insertion and visible reference were preserved.\par
Current source comparison: The visible replacement (7)\_1 is clear; its cancelled predecessor remains unreadable.\par\medskip
\par\noindent\leavevmode\phantomsection\label{GELFAND-0159-U211}\label{GELFAND-P001469:x800-y1425-x1005-y1558}\textbf{\texttt{GELFAND-0159-U211}} (x800-y1425-x1005-y1558).\nopagebreak[4]\par
In item В, after “mobile with” and before λ\_i, a word is crossed out. “coord.” is written above it.\par
Candidates: none recorded. Unresolved in the diplomatic source.\par
Historical source review: Coordinates in the original 2408×3480 image: x=800–1005,y=1425–1558. Checked at 3× and 10×. The corrected entry “coord.” is legible, but the original crossed-out word was not established. Both layers were marked in the body.\par
Current source comparison: The insertion is independently legible. The word it replaces remains hidden.\par\medskip
\par\noindent\leavevmode\phantomsection\label{GELFAND-0159-U212}\label{GELFAND-P001469:x500-y2640-x570-y2745}\textbf{\texttt{GELFAND-0159-U212}} (x500-y2640-x570-y2745).\nopagebreak[4]\par
In (10), below “In the Coulomb syst. [source abbreviation]”, a separate sign immediately before λ\_i is crossed out.\par
Candidates: none recorded. Unresolved in the diplomatic source.\par
Historical source review: Coordinates in the original 2408×3480 image: x=500–570,y=2640–2745. Viewed at 10×. Sloping cancelling strokes obscure the original sign; no reading was established.\par
Current source comparison: The isolated sign remains undeciphered. The following λ\_i is retained without a supplied modifier.\par\medskip
\par\noindent\leavevmode\phantomsection\label{GELFAND-0159-U214}\label{GELFAND-P001469:x438-y860-x2110-y1115}\textbf{\texttt{GELFAND-0159-U214}} (x438-y860-x2110-y1115).\nopagebreak[4]\par
In (9), Delta\_0(b lambda)+a lambda is written as the “general solution of syst. (6)”, without an external factor before Delta\_0. The next line reads “Delta\_0 is one of the functions in example (6)”.\par
Candidates: none recorded. Unresolved in the diplomatic source.\par
Historical source review: Coordinates in the original 2408×3480 image: x=438–2110,y=860–1115. Checked on the whole image; the initial formula was examined again at 5×. The general-solution formula requires mathematical comparison with (6). The reference “example (6)” was also retained literally: examples on the preceding frame P001468 are numbered (7), whereas (6) denotes properties of Delta and Omega. No factors or corrected numbers were added.\par
Current source comparison: The missing external factor and reference (6) are confirmed at 5×/10×. A scaling test distinguishes the addition identity from the extra residue-one normalization needed for (8); it does not claim that a factor is required by every part of (6).\par\medskip
\subsection*{\texttt{GELFAND-P001470}}
\par\noindent\leavevmode\phantomsection\label{GELFAND-0159-U401}\label{GELFAND-P001470:top-left-label}\textbf{\texttt{GELFAND-0159-U401}} (top-left-label).\nopagebreak[4]\par
A densely crossed-out note is visible in the upper-left margin; no original number or sign can confidently be read.\par
Candidates: none recorded. Unresolved in the diplomatic source.\par
Historical source review: Original PNG 2408×3480; region x=130–240,y=80–191 pixels. The whole original was viewed, then the region at up to 8×. SourcePage was left unlabelled; a neighboring page number was not carried onto this frame.\par
Current source comparison: No page number is recovered beneath the cancellation. The source frame remains unlabelled.\par\medskip
\par\noindent\leavevmode\phantomsection\label{GELFAND-0159-U402}\label{GELFAND-P001470:top-line-cancellation}\textbf{\texttt{GELFAND-0159-U402}} (top-line-cancellation).\nopagebreak[4]\par
After “then” in the small top line is one short word covered by several horizontal strokes.\par
Candidates: none recorded. Unresolved in the diplomatic source.\par
Historical source review: Original PNG 2408×3480; region x=1110–1300,y=30–112 pixels. After examination at 10×, the letter remnants did not give a secure reading of the crossed-out word.\par
Current source comparison: The remaining fragments still do not establish the cancelled word.\par\medskip
\par\noindent\leavevmode\phantomsection\label{GELFAND-0159-U403}\label{GELFAND-P001470:top-right-cancellation}\textbf{\texttt{GELFAND-0159-U403}} (top-right-cancellation).\nopagebreak[4]\par
Between “eigenval. [source abbreviation]” and the final H is a crossed-out group of large sloping strokes.\par
Candidates: none recorded. Unresolved in the diplomatic source.\par
Historical source review: Original PNG 2408×3480; region x=1950–2130,y=70–195 pixels. Viewed at 10×; cancellation hides the original signs. The final H on the right is legible and was rendered separately.\par
Current source comparison: The final H remains clear and separate from the unreadable earlier group.\par\medskip
\par\noindent\leavevmode\phantomsection\label{GELFAND-0159-U404}\label{GELFAND-P001470:top-eigenvectors-verb}\textbf{\texttt{GELFAND-0159-U404}} (top-eigenvectors-verb).\nopagebreak[4]\par
In parentheses after “and” and before “eigenvect. [source abbreviation]”, the handwritten word has an initial loop with a descender and joined middle letters.\par
Candidates: знаем [know]. Unresolved in the diplomatic source.\par
Historical source review: Original PNG 2408×3480; region x=388–670,y=166–263 pixels. Neighboring lines and 6× and 10× enlargements were re-examined. “Know” was left as a candidate without reconstructing the word from its meaning.\par
Current source comparison: The candidate remains plausible at 10×, but is still marked uncertain in the body and full ledger.\par\medskip
\par\noindent\leavevmode\phantomsection\label{GELFAND-0159-U405}\label{GELFAND-P001470:lower-number-cancellation}\textbf{\texttt{GELFAND-0159-U405}} (lower-number-cancellation).\nopagebreak[4]\par
After “numbers” and before c\_α, a small group of strokes is crossed out.\par
Candidates: none recorded. Unresolved in the diplomatic source.\par
Historical source review: Original PNG 2408×3480; region x=557–666,y=2350–2476 pixels. Viewed at 10×. The numbers c\_α,λ\_i to the right of the cancellation were rendered separately; the cancelled sign was not reconstructed.\par
Current source comparison: The uncancelled parameter symbols remain independent of the unreadable group.\par\medskip
\par\noindent\leavevmode\phantomsection\label{GELFAND-0159-U406}\label{GELFAND-P001470:lower-equation-cancellation}\textbf{\texttt{GELFAND-0159-U406}} (lower-equation-cancellation).\nopagebreak[4]\par
After “Then (13)” and before the similarity sign is a short crossed-out entry with a descending stroke.\par
Candidates: ур. [eq.; source abbreviation]. Unresolved in the diplomatic source.\par
Historical source review: Original PNG 2408×3480; region x=790–952,y=2650–2780 pixels. Viewed at 10×. The surviving descender permits the abbreviation “ур.”, but dense cancellation prevents asserting the reading.\par
Current source comparison: The descender remains compatible with the candidate abbreviation. The body retains a descriptive cancellation rather than asserting it.\par\medskip
\subsection*{\texttt{GELFAND-P001471}}
\par\noindent\leavevmode\phantomsection\label{GELFAND-0159-U407}\label{GELFAND-P001471:top-initial-cancellation}\textbf{\texttt{GELFAND-0159-U407}} (top-initial-cancellation).\nopagebreak[4]\par
In line (17), a separate densely crossed-out sign precedes “In the case”.\par
Candidates: none recorded. Unresolved in the diplomatic source.\par
Historical source review: Original PNG 2408×3480; region x=310–412,y=66–182 pixels. Viewed at 10×; the original letter or symbol is not confidently distinguishable.\par
Current source comparison: The original sign remains hidden.\par\medskip
\par\noindent\leavevmode\phantomsection\label{GELFAND-0159-U408}\label{GELFAND-P001471:top-line-cancellation}\textbf{\texttt{GELFAND-0159-U408}} (top-line-cancellation).\nopagebreak[4]\par
Beneath “In the case (7)\_1”, a separate short line is crossed out with several horizontal lines. Parts of “lin.”, “fcn.” and the final λ are visible.\par
Candidates: лин. ф-я от λ [a linear function of λ; source abbreviations]. Unresolved in the diplomatic source.\par
Historical source review: Original PNG 2408×3480; region x=445–925,y=195–285 pixels. The line was examined at 6× and 10×. The uncancelled explanation below also contains “Ξ is a linear fcn. of λ”, but that repetition was not treated as proof of hidden letters in the crossed-out line.\par
Current source comparison: The visible parts support the marked candidate, but the neighboring repetition does not resolve the covered letters.\par\medskip
\par\noindent\leavevmode\phantomsection\label{GELFAND-0159-U409}\label{GELFAND-P001471:middle-potential-cancellation}\textbf{\texttt{GELFAND-0159-U409}} (middle-potential-cancellation).\nopagebreak[4]\par
In line (18), a word before “Potent. V” is crossed out. Part of its first sign resembles Cyrillic Е, followed by several visible loops.\par
Candidates: Если [If]. Unresolved in the diplomatic source.\par
Historical source review: Original PNG 2408×3480; region x=372–630,y=1565–1715 pixels. Viewed at 8×. The original word remains uncertain because broad sloping strokes intersect the letters.\par
Current source comparison: The broad cancellation remains across the candidate word. It is not silently inserted into the translated sentence.\par\medskip
\par\noindent\leavevmode\phantomsection\label{GELFAND-0159-U413}\label{GELFAND-P001471:diagram20-faint-strokes}\textbf{\texttt{GELFAND-0159-U413}} (diagram20-faint-strokes).\nopagebreak[4]\par
At the left end of the horizontal line in scheme (20) are thickening and sloping strokes that permit an arrowhead reading. Below this end and the dark point with a minus sign are faint broken strokes down and left. The other bold elements of the scheme are clearly legible.\par
Candidates: A leftward arrow at the left end of the horizontal line. Unresolved in the diplomatic source.\par
Historical source review: Original PNG 2408×3480; region x=340–540,y=2630–2760 pixels. Examined at 6× and 8×. Legible elements were reproduced in TikZ; the ambiguous left end and faint strokes were described by SourceFigure without confidently assigning direction or physical connections. A separate legible arrow on the right-hand part of the line points left.\par
Current source comparison: The faint lower-left remnants remain uncertain at 5×/10×. The existing full native drawing is paired with the wider original crop, including the clear right-hand leftward arrow and M-particles caption; the superseded small detail crop is not the selected comparison.\par\medskip
\par\noindent\leavevmode\phantomsection\label{GELFAND-0159-U414}\label{GELFAND-P001471:margin-below20}\textbf{\texttt{GELFAND-0159-U414}} (margin-below20).\nopagebreak[4]\par
In the left margin below number (20) is a short separate, slightly curved sloping stroke.\par
Candidates: none recorded. Unresolved in the diplomatic source.\par
Historical source review: Original PNG 2408×3480; region x=70–108,y=2750–2799 pixels. Examined at 8×. The separate stroke was recorded without guessing a letter or punctuation mark.\par
Current source comparison: The isolated curve still does not establish a letter or punctuation mark.\par\medskip
\subsection*{\texttt{GELFAND-P001472}}
\par\noindent\leavevmode\phantomsection\label{GELFAND-0159-U410}\label{GELFAND-P001472:upper-atiya-name}\textbf{\texttt{GELFAND-0159-U410}} (upper-atiya-name).\nopagebreak[4]\par
In the subheading “on work [source abbreviation]”, a short surname begins “Ati”; its final joined strokes do not confidently establish whether there is a further i.\par
Candidates: Ати [Ati]; Атии [Atii]. Unresolved in the diplomatic source.\par
Historical source review: Original PNG 2408×3480; region x=627–878,y=470–560 pixels. Re-examined at 8×. Catalog information was not used to expand or correct the visible form of the surname.\par
Current source comparison: The ending remains fused. The English reader retains the literal shortened candidate and does not expand the name from external biographical knowledge.\par\medskip
\par\noindent\leavevmode\phantomsection\label{GELFAND-0159-U411}\label{GELFAND-P001472:middle-flow-cancellation}\textbf{\texttt{GELFAND-0159-U411}} (middle-flow-cancellation).\nopagebreak[4]\par
After “Consider [source abbreviation]”, a word on the main line is crossed out; the clearly legible word “trajectories” is inserted above it.\par
Candidates: none recorded. Unresolved in the diplomatic source.\par
Historical source review: Original PNG 2408×3480; region x=530–830,y=1872–1985 pixels. Viewed at 8×. The upper insertion was rendered in full; the earlier lower entry is hidden by several strokes and was not recovered.\par
Current source comparison: The inserted word remains fully legible while the earlier word remains undeciphered.\par\medskip
\par\noindent\leavevmode\phantomsection\label{GELFAND-0159-U412}\label{GELFAND-P001472:bottom-orientation-cancellation}\textbf{\texttt{GELFAND-0159-U412}} (bottom-orientation-cancellation).\nopagebreak[4]\par
After “orient. [source abbreviation]”, an entry is crossed out with visible lower signs resembling indices 1 and 2; T(x\_1,x\_2) is inserted above it.\par
Candidates: x\_1,x\_2. Unresolved in the diplomatic source.\par
Historical source review: Original PNG 2408×3480; region x=1195–1510,y=3185–3310 pixels. Viewed at 6× and 10×. The upper insertion T(x\_1,x\_2) is legible; the original crossed-out entry was given only as a candidate.\par
Current source comparison: The clear inserted trajectory-space notation is retained. The lower cancelled notation remains conjectural.\par\medskip
\subsection*{\texttt{GELFAND-P001473}}
\par\noindent\leavevmode\phantomsection\label{GELFAND-0159-U601}\label{GELFAND-P001473:mid-cancelled-after-c}\textbf{\texttt{GELFAND-0159-U601}} (mid-cancelled-after-c).\nopagebreak[4]\par
To the right of C in the line about the space of connections is a small densely crossed-out entry. Only letter fragments can be distinguished beneath the strokes.\par
Candidates: none recorded. Unresolved in the diplomatic source.\par
Historical source review: Source 2408×3480 pixels; region x=1390–1570,y=1230–1350, measured from the upper-left corner. Examined on the original, in contextual enlargement 3× and separate enlargement 8×. The pre-cancellation writing is not reliably legible. The insertion “irreducible” and the lower addition “(Hermitian)” were read and preserved.\par
Current source comparison: The cancelled entry remains hidden. Both uncancelled qualifications of connections are retained independently.\par\medskip
\subsection*{\texttt{GELFAND-P001474}}
\par\noindent\leavevmode\phantomsection\label{GELFAND-0159-U602}\label{GELFAND-P001474:upper-class-symbol}\textbf{\texttt{GELFAND-0159-U602}} (upper-class-symbol).\nopagebreak[4]\par
In item 3, after “class for”, the letter before the lower index 2 is heavily retraced, showing an arc and an internal intersection. In the next line after “for”, c₂ is legible.\par
Candidates: c\_2; e\_2. Unresolved in the diplomatic source.\par
Historical source review: Source 2408×3480 pixels; region x=1420–1545,y=515–633. Examined on the original, in context and separately at 8×. The repetition in the next line was not used to silently correct the first sign.\par
Current source comparison: The retraced letter still permits c or e at 10×. The next clear c₂ does not certify the earlier letter.\par\medskip
\par\noindent\leavevmode\phantomsection\label{GELFAND-0159-U603}\label{GELFAND-P001474:mid-cancelled-notation}\textbf{\texttt{GELFAND-0159-U603}} (mid-cancelled-notation).\nopagebreak[4]\par
Before the uncancelled connection sign with a tilde in (9), a symbol is crossed out: a large A-like letter with a lower sign resembling nabla. Several diagonal strokes pass over it.\par
Candidates: A\_\textbackslash{}nabla; A\_\{\textbackslash{}widetilde\textbackslash{}nabla\}. Unresolved in the diplomatic source.\par
Historical source review: Source 2408×3480 pixels; region x=400–615,y=1210–1410. Examined on the original, in context and at 7×. The crossed-out notation was preserved as an uncertain reading; the presence of a tilde in its index cannot be established.\par
Current source comparison: The cancellation still prevents deciding whether the lower nabla originally bore a tilde.\par\medskip
\par\noindent\leavevmode\phantomsection\label{GELFAND-0159-U604}\label{GELFAND-P001474:mid-remark-initials}\textbf{\texttt{GELFAND-0159-U604}} (mid-remark-initials).\nopagebreak[4]\par
In the small addition after “Remark [source abbreviation]”, the initials are И., then a dense sign with two tall strokes and a cross connection, then Г.\par
Candidates: И. М. Г. [I. M. G.]; И. Н. Г. [I. N. G.]. Unresolved in the diplomatic source.\par
Historical source review: Source 2408×3480 pixels; region x=393–610,y=1630–1698. Examined on the original, in contextual enlargement 4× and separately at 8×. The middle letter remains ambiguous.\par
Current source comparison: The middle Cyrillic initial remains ambiguous. Neither a familiar name nor the surrounding attribution is used to select it.\par\medskip
\par\noindent\leavevmode\phantomsection\label{GELFAND-0159-U605}\label{GELFAND-P001474:mid-cancelled-line}\textbf{\texttt{GELFAND-0159-U605}} (mid-cancelled-line).\nopagebreak[4]\par
Between “Look at the critical points [source abbreviations]” and “These are connections with curvature 0”, a short entry is crossed out with several thick nearly horizontal lines. The first sign retains an upper arc and lower loop; fragments of other letters appear above and below the strokes, without forming a reliable reading.\par
Candidates: none recorded. Unresolved in the diplomatic source.\par
Historical source review: Source 2408×3480 pixels; entire entry x=340–1000,y=1760–1870. Examined on the original, in context, at 5× and in two overlapping 10× enlargements. The crossed-out entry was localized and marked in the transcription without reconstructing words from meaning.\par
Current source comparison: The overlapping enlarged details still do not establish the crossed-out words. No transition phrase is supplied.\par\medskip
\par\noindent\leavevmode\phantomsection\label{GELFAND-0159-U606}\label{GELFAND-P001474:lower-equation-label}\textbf{\texttt{GELFAND-0159-U606}} (lower-equation-label).\nopagebreak[4]\par
To the left of the formula for ∇ad, the parenthesized number has a legible first 1 and a heavily corrected second digit.\par
Candidates: 11. Unresolved in the diplomatic source.\par
Historical source review: Source 2408×3480 pixels; region x=145–325,y=2455–2610. Examined on the original, in context and at 8×. The candidate 11 was preserved as an uncertain reading; the original corrected digit was not recovered.\par
Current source comparison: The final candidate 11 remains marked uncertain; the overwritten earlier digit is not reconstructed.\par\medskip
\par\noindent\leavevmode\phantomsection\label{GELFAND-0159-U607}\label{GELFAND-P001474:lower-m-x-notation}\textbf{\texttt{GELFAND-0159-U607}} (lower-m-x-notation).\nopagebreak[4]\par
Formula (10) directly reads Ω¹(M). The next formula, with corrected number tentatively (11), directly reads Ω⁰(X) and Ω¹(X). The numerator and denominator of (13) likewise read Ω¹(X) and Ω⁰(X), respectively.\par
Candidates: Ω¹(M) in (10); Ω⁰(X), Ω¹(X) in (11); Ω¹(X), Ω⁰(X) in (13).. Unresolved in the diplomatic source.\par
Historical source review: Notation checked against the 2408×3480 source: formula (10), x=780–1870,y=2260–2445; formula (11), x=440–1850,y=2460–2650; formula (13), x=365–1370,y=2880–3160. All three regions were directly re-examined at 3×. M and X were read confidently. This is a record of external notation checking: X was retained verbatim; no new equality or substantive correction was introduced into the formulas.\par
Current source comparison: The additional 5×/10× details confirm M in (10) and X in both positions in (13). The notation is legible, while compatibility of the spaces is a separate mathematical question.\par\medskip
